% !TEX root = ../main.tex

%----------------------------------------------------------------------------------------
% CHAPTER 1
%----------------------------------------------------------------------------------------

\chapter{Einleitung}
\label{Chapter1}

%----------------------------------------------------------------------------------------

\section{Ausgangslage und Zielsetzung}
\label{sec:ausgangslage}

Das Swiss Household Panel (SHP) ist eine jährliche Panelbefragung, die seit 1999
durchgeführt wird. Koordiniert wird sie vom Schweizerischen Kompetenzzentrum für
Sozialwissenschaften FORS in Lausanne. Erhoben werden Lebensbedingungen, soziale
Einbettung, politische Einstellungen und subjektive Bewertungen. Dieselben Personen
werden über mehrere Wellen hinweg befragt \cite{FORS2024}. Die Stichprobe beruht auf
einer Zufallsauswahl von Privathaushalten in der ganzen Schweiz. Um selektive
Panelattrition auszugleichen, wird sie in regelmässigen Abständen durch
Auffrischungsstichproben ergänzt \cite{Voorpostel2018}. Weil die Befragung
längsschnittlich angelegt ist, lassen sich nicht nur Querschnittsvergleiche zwischen
Gruppen, sondern auch individuelle Veränderungen über die Zeit beobachten. Das SHP
gehört damit zu den zentralen Datenquellen für die Analyse sozialer und politischer
Entwicklungen in der Schweiz.

Die Fragestellung dieses Projekts lautet: Wie hat sich das Vertrauen der Schweizer Wohnbevölkerung in den Bundesrat seit 1999 entwickelt, und gibt es
zwischen verschiedenen Untergruppen systematische Unterschiede? Die Zielvariable
PP04\footnote{Ist die Variable aus dem Datensatz. Siehe Anhang \ref{AppendixA}.} misst dieses Vertrauen auf einer elfstufigen Skala von 0 (kein Vertrauen)
bis 10 (volles Vertrauen). Sie wird im SHP in unregelmässigen Abständen erhoben,
erstmals 1999 und zuletzt 2023.

Ausgangspunkt der Auswertung sind die SHP-Rohdaten. Daraus wird ein bereinigter
Paneldatensatz auf Personen-Jahr-Ebene aufgebaut. Er enthält alle in der Aufgabenstellung
vorgegebenen Variablen sowie eine zusätzliche Einflussgrösse, die vom Projektteam
gewählt wurde.

Ein besonderes Augenmerk liegt auf der Reproduzierbarkeit. Alle Verarbeitungsschritte,
von der ursprünglichen ZIP-Datei bis zum fertigen Analysedatensatz, sind in R
programmiert und im Projektverzeichnis abgelegt.

%----------------------------------------------------------------------------------------

\section{Hypothese}
\label{sec:hypothese}

Als zusätzliche Einflussgrösse wählt das Projektteam den Zivilstand.
Der Zivilstand prägt die soziale Einbettung einer Person und steht damit plausibel in
Zusammenhang mit dem Vertrauen in staatliche Institutionen. In der Literatur weisen
verheiratete Personen im Schnitt höhere Werte für generalisiertes Vertrauen auf, während
Trennungen mit tieferem institutionellen Vertrauen assoziiert sind.

Daraus leitet sich folgende Hypothese ab:

\begin{quote}
Verheiratete Personen weisen im Durchschnitt ein höheres Vertrauen in den Bundesrat
auf als ledige, getrennte oder geschiedene Personen.
\end{quote}

Die Prüfung erfolgt in Kapitel~\ref{Chapter3} deskriptiv über gruppenspezifische
Mittelwerte und Mediane von PP04 sowie über die zeitliche Entwicklung je
Zivilstandsgruppe.